\chapter{Detalle de la evaluación de ML-2}
\label{anx:ml2}

Este anexo completa la Sección~\ref{sec:ml2} con la construcción del conjunto de datos, la
configuración de los entrenamientos, los resultados completos y la revisión manual de falsos
positivos.

\section{Conjunto de datos}

Cada mención de técnica marcada en una frase de AnnoCTR se empareja, por orden, con su enlace a
ATT\&CK del mismo informe. Casan 3.374 de 3.375 menciones. Los identificadores se normalizan a la
técnica padre vigente de la v19.2: se traducen 30 revocados y se descartan 18 de ATT\&CK Mobile o
desconocidos. Los 280 informes sin anotar de AnnoCTR no se usan.

\begin{table}[H]
\centering
\begin{tabular}{|l|r|r|r|}
\hline
\textbf{Partición} & \textbf{Informes} & \textbf{Frases} & \textbf{Frases con técnica} \\
\hline
Entrenamiento & 70 & 7.570 & 1.499 \\
Desarrollo & 16 & 1.550 & 335 \\
Prueba & 34 & 3.059 & 754 \\
\hline
\end{tabular}
\caption{Particiones de AnnoCTR, tal como se publican.}
\end{table}

\section{Configuración de los entrenamientos}

Los dos modelos se afinan con la misma configuración: clasificación multietiqueta con pérdida de
entropía cruzada binaria, AdamW (tasa de aprendizaje $3\cdot10^{-5}$, decaimiento 0,01, 10\,\% de
calentamiento lineal), lotes de 32 frases de hasta 128 \textit{tokens}, 8 épocas, precisión mixta
\texttt{bfloat16} y recorte del gradiente a 1. En cada época se calcula en desarrollo el F1 por
frase para 18 umbrales (de 0,05 a 0,90) y se guarda la mejor combinación de época y umbral. Las
épocas elegidas van de 4 a 7 y los umbrales de 0,10 a 0,25. El entorno es Python 3.12 con PyTorch
2.11 (CUDA 12.8) y Transformers 5.20 sobre una RTX~5070 de 8\,GB.

\section{Resultados completos}

\begin{table}[H]
\centering
\small
\resizebox{\textwidth}{!}{%
\begin{tabular}{|l|c|c|c|c|c|}
\hline
\textbf{Sistema (prueba)} & \textbf{P frase} & \textbf{R frase} & \textbf{F1 frase} & \textbf{F1 macro} & \textbf{F1 informe} \\
\hline
Diccionario sin poda & 0,203 & 0,219 & 0,211 & 0,325 & 0,391 \\
Diccionario podado & 0,720 & 0,214 & 0,330 & 0,349 & 0,490 \\
ModernBERT-base & 0,568 & 0,480 & 0,519 $\pm$ 0,010 & 0,262 & 0,620 $\pm$ 0,041 \\
SecureBERT 2.0 & 0,612 & 0,506 & 0,553 $\pm$ 0,019 & 0,301 & 0,670 $\pm$ 0,026 \\
\hline
\end{tabular}%
}
\caption{Resultados en prueba, con media y desviación de tres semillas.}
\end{table}

\begin{table}[H]
\centering
\small
\resizebox{\textwidth}{!}{%
\begin{tabular}{|l|c|c|}
\hline
\textbf{Comparación (prueba)} & \textbf{Dif. F1 frase [IC 95\,\%]} & \textbf{Dif. F1 informe [IC 95\,\%]} \\
\hline
SecureBERT 2.0 $-$ ModernBERT & +0,034 [+0,011, +0,058] & +0,050 [+0,027, +0,073] \\
SecureBERT 2.0 $-$ diccionario & +0,223 [+0,132, +0,329] & +0,179 [+0,071, +0,317] \\
ModernBERT $-$ diccionario & +0,189 [+0,092, +0,295] & +0,130 [+0,020, +0,267] \\
\hline
\end{tabular}%
}
\caption{Diferencias con remuestreo emparejado por informe (10.000 muestras).}
\end{table}

\section{Decisión de la combinación}

\textbf{Reglas, fijadas antes de ejecutar la validación cruzada:}
\begin{itemize}
    \item cinco bloques por informe sobre los 86 de entrenamiento y desarrollo;
    \item SecureBERT~2.0 con la época y el umbral medianos de los entrenamientos anteriores (7 y 0,20), sin elegir nada dentro de cada vuelta;
    \item la poda del diccionario y el corte de «técnica rara» se calculan solo con los bloques de entrenamiento de cada vuelta;
    \item la combinación se adopta solo si mejora y el remuestreo da menos de un 5\,\% de diferencias nulas o negativas.
\end{itemize}

\begin{table}[H]
\centering
\small
\resizebox{\textwidth}{!}{%
\begin{tabular}{|l|c|c|c|c|c|}
\hline
\textbf{Variante (86 informes)} & \textbf{F1 frase} & \textbf{F1 macro} & \textbf{R explícitas} & \textbf{R implícitas} & \textbf{F1 informe} \\
\hline
Diccionario & 0,249 & 0,278 & 0,327 & 0,096 & 0,479 \\
Modelo & 0,481 & 0,206 & 0,600 & 0,360 & 0,620 \\
Modelo $\cup$ diccionario en técnicas con $<$20 frases & 0,509 & 0,373 & 0,717 & 0,372 & 0,690 \\
\textbf{Modelo $\cup$ diccionario} & \textbf{0,520} & \textbf{0,382} & \textbf{0,771} & \textbf{0,387} & \textbf{0,695} \\
\hline
\end{tabular}%
}
\caption{Validación cruzada por informe. Unión frente a modelo: +0,040 de F1 por frase
[+0,020, +0,060] y +0,076 por informe [+0,042, +0,109].}
\end{table}

En desarrollo, el diccionario saca un F1 por frase de 0,12, frente a 0,33 en prueba y 0,25 en la
validación cruzada. Desarrollo tiene un 24\,\% de menciones explícitas y prueba un 34\,\%.

\section{Errores por frecuencia}

\begin{table}[H]
\centering
\begin{tabular}{|l|r|c|c|}
\hline
\textbf{Frases de entrenamiento} & \textbf{Técnicas} & \textbf{F1 medio, modelo} & \textbf{F1 medio, diccionario} \\
\hline
5--9 & 16 & 0,120 & 0,349 \\
10--19 & 13 & 0,247 & 0,472 \\
20--49 & 13 & 0,466 & 0,294 \\
50 o más & 7 & 0,558 & 0,222 \\
\hline
\end{tabular}
\caption{F1 por técnica según su frecuencia en entrenamiento (SecureBERT~2.0, tres semillas).}
\end{table}

Las confusiones más frecuentes del modelo enfrentan técnicas vecinas de una misma cadena:
intérprete de órdenes con \textit{phishing}, transferencia de herramientas con intérprete,
ejecución por el usuario con \textit{phishing}, y codificación de datos con ofuscación o con canal
cifrado.

\section{Revisión manual de falsos positivos}

Se revisó una muestra aleatoria fija de 60 de los 270 falsos positivos del sistema adoptado en
prueba (una semilla), en tres clases:
\begin{itemize}
    \item \textbf{H, hueco de anotación:} la frase describe o nombra la técnica y AnnoCTR no la anotó;
    \item \textbf{D, dudoso:} solo es correcta con el contexto del informe, o la técnica es vecina de la anotada;
    \item \textbf{E, error:} la frase no describe la técnica.
\end{itemize}

\begin{table}[H]
\centering
\begin{tabular}{|l|r|r|r|}
\hline
\textbf{Clase} & \textbf{Total} & \textbf{Del modelo} & \textbf{Del diccionario} \\
\hline
H & 25 (42\,\%, IC 95\,\% 30--54\,\%) & 16 & 9 \\
D & 10 (17\,\%) & 10 & 0 \\
E & 25 (42\,\%) & 18 & 7 \\
\hline
\end{tabular}
\caption{Clasificación de la muestra de falsos positivos.}
\end{table}

Contando la clase H como acierto, la precisión por frase pasa de 0,61 a 0,77 (0,73--0,82); con la
clase D, a 0,84. La exhaustividad no cambia, porque la revisión no busca técnicas que el sistema
tampoco ve. Los huecos típicos son filas de tablas de ATT\&CK y técnicas nombradas en listas,
descargas de herramientas, activación de macros y órdenes de PowerShell o \texttt{rundll32}
literales. Los errores del diccionario vienen de nombres genéricos de subtécnica que
sobrevivieron a la poda, como «Web Services» o «Exploits». Se informa como oficial la precisión
medida y la corregida como estimación.

\section{Perfil de APT29}

De los seis informes de Microsoft se descartan las secciones de mitigación, detección, consultas
de búsqueda, productos e indicadores, y todo lo que sigue a «Related posts». Quedan 683 frases de
análisis. Las tablas «MITRE ATT\&CK techniques observed» del propio editor se excluyen del perfil
principal, porque alimentan en parte la ficha del grupo.

\begin{table}[H]
\centering
\begin{tabular}{|l|r|c|c|c|}
\hline
\textbf{Sistema} & \textbf{Técnicas} & \textbf{Precisión} & \textbf{Exhaustividad} & \textbf{F1} \\
\hline
Diccionario & 12 & 0,92 & 0,33 & 0,49 \\
SecureBERT 2.0 & 16 & 0,81 & 0,39 & 0,53 $\pm$ 0,09 \\
\textbf{SecureBERT 2.0 $\cup$ diccionario} & 24 & 0,83 & 0,60 & 0,69 $\pm$ 0,05 \\
\hline
\end{tabular}
\caption{Perfil de APT29 frente a las 33 técnicas de G0016 y sus campañas que caben en el espacio
de etiquetas (técnicas con al menos una frase).}
\end{table}

El modelo no extrae el compromiso de la cadena de suministro (\texttt{T1195}) del informe de
SolarWinds, pese a tener 20 frases de entrenamiento; el diccionario sí. En el paquete STIX, cada
técnica del perfil es una relación \texttt{uses} de G0016 con su confianza, el número de frases y
de informes que la respaldan, el método que la extrajo y una frase de evidencia. La confianza es
la probabilidad máxima del modelo; cuando la técnica solo la extrae el diccionario, se usa 61, su
precisión en la validación cruzada.

\section{Aumento con procedimientos de ATT\&CK}

Las 343 frases de procedimiento (189 de \texttt{T1003}, 74 de \texttt{T1548} y 80 de
\texttt{T1569}) salen de las relaciones de uso de ATT\&CK v19.2, sin las de APT29 ni sus campañas.
Solo entran en el entrenamiento de cada vuelta y nunca se predicen. Los dos brazos usan los mismos
cinco bloques, semillas, época, umbral y diccionario podado, así que solo cambia el modelo. El
brazo sin aumento reproduce exactamente la validación cruzada adoptada.

\begin{table}[H]
\centering
\begin{tabular}{|l|c|c|c|c|c|}
\hline
\textbf{Sistema (86 informes, 59 técnicas)} & \textbf{F1 frase} & \textbf{P} & \textbf{R} & \textbf{F1 macro} & \textbf{F1 informe} \\
\hline
Modelo & 0,481 & 0,585 & 0,408 & 0,206 & 0,620 \\
Modelo $\cup$ diccionario & 0,520 & 0,569 & 0,480 & 0,382 & 0,695 \\
Modelo aumentado & 0,470 & 0,615 & 0,380 & 0,199 & 0,607 \\
\textbf{Modelo aumentado $\cup$ diccionario} & \textbf{0,518} & \textbf{0,593} & \textbf{0,460} & \textbf{0,382} & \textbf{0,700} \\
\hline
\end{tabular}
\caption{Validación cruzada del aumento. Diferencia de las uniones: $-0{,}002$ por frase
[$-0{,}018$, $+0{,}013$] y $+0{,}005$ por informe [$-0{,}012$, $+0{,}022$].}
\end{table}

El modelo aumentado gana precisión y pierde algo de exhaustividad, sin diferencia significativa
($-0{,}011$ por frase), y el diccionario lo compensa en la unión. El perfil de APT29 regenerado con el
sistema aumentado tiene 22 técnicas: entra \texttt{T1548} y sale \texttt{T1055}, y el F1 frente a
G0016 y sus campañas pasa de 0,655 a 0,667, con precisión 0,86.

\section{Perfiles de los otros actores}

\begin{table}[H]
\centering
\small
\begin{tabular}{|l|l|r|r|r|c|c|c|}
\hline
\textbf{Nombre de Microsoft} & \textbf{Grupo de ATT\&CK} & \textbf{Técnicas} & \textbf{Ficha} & \textbf{Aciertos} & \textbf{P} & \textbf{R} & \textbf{F1} \\
\hline
Forest Blizzard & G0007 APT28 & 8 & 33 & 5 & 0,63 & 0,15 & 0,24 \\
Mint Sandstorm & G0059 Magic Hound & 8 & 35 & 6 & 0,75 & 0,17 & 0,28 \\
Peach Sandstorm & G0064 APT33 & 4 & 18 & 2 & 0,50 & 0,11 & 0,18 \\
Octo Tempest & G1015 Scattered Spider & 5 & 26 & 4 & 0,80 & 0,15 & 0,26 \\
Volt Typhoon & G1017 Volt Typhoon & 9 & 35 & 7 & 0,78 & 0,20 & 0,32 \\
Storm-1811 & G1046 Storm-1811 & 8 & 14 & 5 & 0,63 & 0,36 & 0,45 \\
Storm-0501 & G1053 Storm-0501 & 9 & 18 & 9 & 1,00 & 0,50 & 0,67 \\
\hline
\end{tabular}
\caption{Perfil extraído de un informe frente a la ficha del grupo con sus campañas, dentro del
espacio de 62 técnicas.}
\end{table}

Los actores menos documentados tienen fichas más cortas, y un solo informe cubre una parte mayor
de ellas. Storm-0558, que no tiene grupo en ATT\&CK, y el informe sobre ransomware operado por
personas, que trata de varios actores, entran en el grafo solo como informes. Aquí se repite la
circularidad de APT29: las fichas de ATT\&CK citan en parte estos mismos informes.
